\section{Structural identity and the magnitude of mass}
\label{vi:sec:introduccion}

Determining a particle's mass requires specifying which object persists
and which operation converts its structure into a magnitude. This work
constructs that connection: compatible extensions and their records
produce fibres of states, on which a positive mass operator is defined.
Its electronic specialisation brings together two coincident readers in
one scalar evaluation. The thesis is that identity, the mass equation,
and the evaluated value belong to a common mathematical chain; physical
comparison subsequently examines the realisation of that chain in a
measurement chart. The equation's arguments and coefficients thus retain
a structural provenance preceding comparison.

Trajectory dynamics incorporates history into that magnitude: it determines
the exact cellular support of the central trajectory, reconstructs two
memories that a counting statistic cannot distinguish, diagonalises the return
reader by finite characters, and characterises mass conservation by the
commutator between the mass operator and a TPK transport. These four operations
act on the same state and do not replace the atlas, the electron equation, or
their prior definitions and proofs.

A formal mathematical kernel called Triadic Modular Holography,
henceforth HMT, is introduced through three objects. \emph{All Possible
Paths}, henceforth APP, is the arithmetic path structure on
\((\mathbb Z/9\mathbb Z)^2\). Adjacency is given by the additive Cayley
graph with cardinal generators \(\pm(1,0)\) and \(\pm(0,1)\); the
periodic cellular realisation of that support is toroidal. Its
operations of additive accumulation and multiplicative composition are
distinguished from the evaluations determining their coordinates. TRIT
is the unit of oriented signature \(\{-1,0,+1\}\): it determines the
local regime and orientation of an operation. The \emph{Topological Phase
Kernel}, henceforth TPK, constitutes the kernel's dynamics; it composes
selection, transport, memory updating, and return on the enriched state.
Dimensional Mechanics develops realisations of that structure on
magnitude lines and in physical state spaces. The exposition constructs
these objects from their rules before using the internally obtained coordinates
in the mass operators.

The development obtains definite results at each of these levels. The
composition of \(108\) TPK updates produces a terminal state
that preserves the incidence ledger, its two transverse channels, and the
dodecaphase signature; the orbital transform reversibly reconstructs the record
\(K\). This genealogy supports the construction of the
\(81/56/13/324\) atlas, the positive mass operator of each fibre, the
scalar exception of the central electron, and the composition rules preceding
evaluation. The subsequent comparison keeps central values,
intervals, bounds, pole masses, running masses, and widths separate. The contribution
therefore consists in an effective chain linking the persistent object,
its route memory, the operator acting on its fibre, and the precise type of
the evaluated magnitude.

In that same causal order, the complete census
\(104{,}976\to468\to243\subset729\) precedes every regional selection. The
bijection \(\beta:\mathbb F_3^6\to E_9^3\) identifies the ambient space of
\(729\) positions, whereas the threshold difference produces the shell
\(271=243+27+1\) and the exact completion
\(1000=729+243+27+1\). The uniform measure of the completion is compatible
with the dimensional projections. These operations construct the ambient
space and its boundary; subsequent selection retains its own domain.

\subsection{The structure preceding the numerical value}

Mass is not the only information carried by a particle. The
support of persistence, its orientations, conjugation,
representations, and composition rules also enter. In the HMT approach,
these data are not appended at the end to an isolated numerical value: they belong to the
construction on which the mass functional acts. Two histories can
share a scalar reading while preserving different information in their
sheets or memory. Equality of values must therefore be distinguished
from equality of structures.

This distinction is particularly clear in composition. A table can
add coordinates that have already been evaluated; a dynamical construction must first determine whether two
histories can compose their phase, boundary, orientations, and
predecessors. Records are composed at that level. The signature is extracted
afterwards, and the character evaluates that signature. The distinction between these
operations provides the mathematical setting for studying composite states,
bindings, and resonances.

The question of why a magnitude has the value it does is thus linked to
a more precise question: which elements of the construction determine its value, and
what information remains available after its evaluation? The article
develops this question from the internal rules of the corpus. The
subsequent physical comparison examines their realisations without intervening in
the selection of the coefficients subjected to comparison.

\subsection{Hierarchy of the exposition}

The formal kernel is reproduced in common with the preceding articles
in the series. This repetition serves self-containment:
each text must be readable without access to the others. The specific
extensions of the present work follow the kernel and are
identified by the objects they construct.

The generated coordinates, the fine-structure constant, and the reduced
Planck constant enter when the mass equations
require them. The angular pair precedes the channels and their moments;
the records precede the signature; the fibres precede the operators
acting on them. The central electron is preserved as a structural
construction in its own right, even when its magnitude is expressed on a shared
chronological scale.

Within this hierarchy, selection, transport, and updating act on the same
enriched state. The composition of \(108\) steps produces the terminal
state; its projection onto the dodecaphase ledger, evaluation of the
\(90/120\) channels, the signed extraction \(U_{12}^{\rm sgn}\), the
orbital transform, and integral reconstruction establish the chain
\[
 X_{\rm seed}\longmapsto X_{\rm term}\longmapsto
 C_{\rm term}\longmapsto U_{\rm term}\longmapsto K.
\]
The record, its \(12\) components, its quotients, and its memory arise
from that evolution. The channel, signature, and period charts preserve
within one domain the sheet, orientation, carry, and boundary required by
mass evaluation. Composition~\eqref{vi:dep:cadena-terminal-completa} brings
together the producer and its reconstruction, with the adjacent proofs of
updating and inversion; identity~\eqref{vi:dep:K-reversible} specifies the
reversible transform and its integral sublattice.

The definition of a particle prepares the classification. Conjugation
and holonomy, the atlas, the internal observables, and mass evaluation
are then developed. The organisation by collections uses
that architecture, so that each sector is presented with its support,
operation, and reading. Kepler--Sommerfeld is situated within the orbital
development of action, where the maps and conditions
justifying the relation can be specified; it is not introduced as a
historical resemblance devoid of operational content.

\subsection{Comparison as the result of an explicit chain}

Comparison requires more than two columns of numbers. One must know what
has been evaluated, in which unit, and under which definition of the observable. A
pole mass, a running mass at a stated scale, and an
experimental bound are not interchangeable. A width expresses
dynamic information different from the position of a mass. The size of the inventory
is useful when it preserves these distinctions, since it makes it possible to examine
complete collections rather than only favourable examples.

For this reason, the presentation brings together each output, its provenance,
and the conditions of its realisation. An internal identity, the
evaluation of a specified signature, and the identification of that signature
with a species are related assertions with proofs
of their own. Presenting them consecutively shows what the
model determines at each level and the specific content of the comparison.

The argumentative unity of the article is, ultimately, the continuity between
persistence, classification, operation, and measurement. Mass is an output
of this continuity; the structure producing it remains available
for understanding its relations with the other readings of the particle.

% BEGIN SINTESIS_ADITIVA_20260922
\par
Section~\ref{delta22:vi:composicion} incorporates structural length before
the clock: marked incidence produces \(N=5760\),
\(r_N=5759/23040\), and
\(L_\alpha=\alpha_{\rm HMT}^{16}r_NL_\star\)
in~\eqref{delta22:vi:Lalpha}; only then are \(t_0\),
\(R_{108}=L_\alpha\), \(Q_L\), \(m_{\rm clk}\), and
\(b_e(Q_L)\) fixed. The cotransformation law
\eqref{delta22:vi:be-cotransformacion} shows why the electronic factor
cannot be frozen when the longitudinal or action section changes. The
operator relations \eqref{delta22:vi:R01} and
\eqref{delta22:vi:R02} eliminate action length and determine the unique
gravitational coefficient~\eqref{delta22:vi:G-unica}; the circular
formula~\eqref{delta22:vi:G-circular} follows afterwards. Mass, its route
character, and its tables remain causally prior to that radial
identification.
% END SINTESIS_ADITIVA_20260922
